\chapter*{不定积分每日一题}

\refer{《积分的方法与技巧》金玉明\ 顾新身\ 毛瑞庭\ /\ 编著}

\exercise{求积分$\int{\frac{1}{\sqrt{x}\qty(1+\sqrt[3]{x})\d x}}$}

\teacher{陈通涨}

\postdate{2022年12月5日}

\result{6\qty(\sqrt[6]{x}-\arctan{\sqrt[6]{x}})+C}

\analysis{这是一道非常经典的换元积分法例题.考虑换元后将原积分转化为更熟悉的有理分式积分，再进行求解.}

\solution{
    在本题中，为了将式子化为有理分式，我们可设$t=\sqrt[6]{x}$，则有
    \formula{
        \addtolength{\jot}{-1em}
        \begin{aligned}
            \int{\frac{6t^5}{t^3\qty(1+t^2)}\d t}&=6\int{\frac{t^2}{1+t^2}\d x}\\
            &=6\int{1-\frac{1}{1+t^2}\d t}\\
            &=6\qty(t-\arctan{t})+C\\
            &=6\qty(\sqrt[6]{x}-\arctan{\sqrt[6]{x}})+C
        \end{aligned}
    }
}
    上面的题目比较简单，我们接下来尝试一道运用类似换元法的题.
        \formula{\int{\frac{1}{\sqrt[3]{\qty(x-1)\qty(x+1)^2}}\d x}}
    \analysis{
        我们也可以考虑换元法变成有理分式，但与前面相比得先变形一下
        \formula{\frac{1}{\qty(x-1)^{\frac{1}{3}}\qty(x+1)^{\frac{2}{3}}}=\frac{\qty(x+1)^{\frac{1}{3}}}{\qty(x-1)^{\frac{1}{3}}}\cdot\frac{1}{x+1}=\qty(\frac{x+1}{x-1})^{\frac{1}{3}}\cdot\frac{1}{x+1}}
        这时再考虑整体换元，令$\qty(\frac{x+1}{x-1})^{\frac{1}{3}}=t$，则$x=\frac{1+t^3}{t^3-1}$，后面就可以操作了.解法如下：
        \formula{
            \addtolength{\jot}{-1em}
            \begin{aligned}
                \int{\frac{1}{\sqrt[3]{\qty(x-1)\qty(x+1)^2}}\d x}=&\int{t\cdot\frac{t^3-1}{2t^3}\d\qty(\frac{1+t^3}{t^3-1})}\\
                =&\int{\frac{-3}{t^3-1}\d t}\\
                =&\int{\frac{-3}{\qty(t-1)\qty(t^2+t+1)}\d t}\\
                =&-\int{\frac{1}{t-1}\d t+\int{\frac{1+2t}{t^2+t+1}\d t}}\\
                =&-\ln{\abs{t^2+t+1}}+\frac{1}{2}\ln{\abs{t^2+t+1}}+\sqrt{3}\arctan{\frac{2t+1}{\sqrt{3}}}+C\\
                =&-\ln{\abs{\sqrt[3]{\frac{x+1}{x-1}}-1}}+\frac{1}{2}\ln{\abs{\sqrt[3]{\qty(\frac{x+1}{x-1})^2}+\sqrt[3]{\frac{x+1}{x-1}}+1}}\\
                &+\sqrt{3}\arctan{\qty(\frac{2}{\sqrt{3}}\sqrt[3]{\frac{x+1}{x-1}})}+C
            \end{aligned}
        }
    }